\section{Consequences and limits}\label{sec:discussion}
The joint requirement identifies a specific role for global advice. At a fixed load, $1/(1-\rho)$ is the infimum maximum-flow factor compatible with the proved tail-order objective. A finite message chooses among finitely many such factors, and the best inflation over a bounded load envelope is governed by logarithmic covering. The geometric bins are not themselves a new scheduling principle. They become meaningful because a universal capacity obstruction and a nonclairvoyant construction meet on opposite sides of the same strict boundary.

\subsection{What the boundary says about information}
The advice message supplies coarse environmental information, not clairvoyance about jobs.  Once a code is selected, guarded sharing uses only release order and the active population.  Its deterministic certificate is valid on every finite input, including inputs that have zero probability under the advertised environment.  Correct advice is used only to place the residual PS rate above the stochastic load.  This division of labor differs from a predictor that supplies per-job sizes, ranks, or future arrivals, and from an online-advice tape constructed after seeing the realized sequence.

The lower bound shows that this coarse information cannot compensate for a factor below the load boundary.  It permits the competing policy to observe sizes of arrived jobs, to maintain arbitrary within-cycle history, and to depend on the selected environment code.  What it forbids is future realization information and cross-cycle state outside the declared reset.  Thus the obstruction is not caused by nonclairvoyance of the positive rule.  It is caused by the amount of capacity a finite-prefix maximum-flow deadline forces onto a large root.

At the covering infimum, $\log\Gamma_M^*(u,v)=\log(v/u)/M$: the covered coordinate is $\log\theta$, not $\log M$.  This does not make all distributional information one-dimensional.  The tail constant in Theorem~\ref{thm:guard-tail} also depends on the regular-variation index through the analysis, and finer objectives could depend on the slowly varying factor or on the entire law.  Our result says that, for the stated order objective and construction, one load bin suffices; it does not say that load is a sufficient statistic for every scheduling goal.

\subsection{Relation to richer scheduling information}
Service-time information can be used in many other ways.  Two-class and rank-based policies interpolate between blind and fully size-aware scheduling~\cite{chen2021,scully2018soap}; Gittins and related indices optimize stochastic mean objectives under distributional information~\cite{gittins2011,aalto2009,scully2020gittins}; and large-deviation priority rules optimize different rare-event criteria~\cite{stolyar2001}.  These mechanisms may improve constants or objectives that guarded sharing does not target.  None directly supplies the exact maximum-flow certificate proved by the oldest-job reserve, because that certificate protects every released prefix rather than an average or asymptotic state.

Conversely, the reserve is not proposed as universally efficient.  It deliberately spends a fixed rate on the oldest job even when a size-aware rule would prefer another job.  The point is analytical compatibility: the reserve gives a pathwise prefix deadline, while the shared rate gives every job a stable comparison queue.  Any alternative policy that improves one side must still expose an interface strong enough to prove the other.  A promising route is to replace the constant reserve by a state-dependent certificate whose cumulative service to every prefix can be lower-bounded, but the tail comparison would then have to tolerate a time-varying residual capacity.

\subsection{Which hypotheses are structural}
Several distinctions limit how the theorem should be used. The deterministic guarantee concerns the maximum response, not total response or a competitive fairness criterion. The stochastic guarantee concerns a typical arrival in an M/G/1 queue, not worst-case release sequences. Its service laws have regularly varying tails with index below $-1$, and the policy class has an empty-state reset and no future-horizon input. A fixed code remains maximum-flow competitive under incorrect advice, but it need not retain the tail guarantee. Removing any of these restrictions is a separate theorem obligation.

The Poisson assumption is more than a convenient arrival rate.  It supplies the stationary-arrival sampling and invariant PS environment used by the upper bound, and the independent future marked process used by the lower bound.  Renewal arrivals may preserve the fluid event but change what an arrival sees and how a permanent tagged customer is initialized.  A non-Poisson extension therefore needs a Palm version of both halves, not only a renewal strong law.

The single-server assumption is equally structural.  On parallel servers, unfinished work no longer determines maximum flow: placement and fragmentation matter, and an oldest job may be blocked behind work assigned elsewhere.  Multiserver SRPT and Gittins results show that even stochastic response analysis requires new coupling ideas~\cite{grosof2018,scully2020gittins,hong2024}.  An advice boundary in that setting would need a new adversarial benchmark and likely more than one scalar environmental coordinate.  The present formula $(v/u)^{1/M}$ should not be extrapolated to it.

The construction deliberately uses a stronger tail-analysis comparison than its rate rule needs to implement. Every job receives enough shared service to be compared with a stable PS queue, while an oldest-job reservation protects all prefixes. The permanent-customer proof shows why a high conditional response moment can be useful even when the corresponding service moment is infinite: the moment is first taken over a light-tailed customer count, conditional on service size, and only then integrated through a truncated moment. Interchanging these steps would demand a nonexistent high service moment and lose the regularly varying cases of interest.

The converse has a complementary structure. It counts old bounded jobs left after a forced large-job completion. It does not equate backlog volume with queue length, and it does not count overlapping large-job events throughout a stationary history. Restricting the trigger to one root per regenerative cycle makes the reward integrable and the conversion to a Palm arrival probability explicit. This is the step that yields the factor $x$ separating the two tail orders.

\subsection{Open analytical boundaries}
The equality boundary $C=1/(1-\rho)$ remains unclassified. Both present proofs lose their strict capacity margin there, for different reasons. The exact finite-advice statement is therefore an infimum, including at public bin endpoints. Resolving equality may require a policy whose spare capacity vanishes with state or age, together with a tail proof that does not compare against a strictly stable slower queue. A converse would need to extract an order-$x$ cohort without choosing a truncation whose load exceeds $1-1/C$, which is impossible at equality by monotone convergence. Neither direction follows by taking a limit in the existing theorems.

An optimal leading tail constant is also not determined. The supplied upper constants establish order and quantify sufficient slack, but the permanent-customer enlargement and Markov inequality are not expected to be sharp. Classical PS asymptotics suggest that stronger distributional assumptions could improve constants~\cite{yashkov1987,zwartboxma2000,borst2006}. A matching lower constant for the joint objective would require optimizing over all eligible policies, not only analyzing guarded sharing.

Randomization is another genuine boundary.  A pathwise competitive factor for every random seed would reduce to the deterministic theorem seed by seed.  An expected factor, or a high-probability factor against an oblivious adversary, has different quantifiers and does not force the large root to finish on each fluid-event path.  Establishing a randomized analogue would require stating the adversary's access to the seed and replacing the deterministic cohort argument accordingly.

There is a further difference between global advice and a learned scheduler. A message may be selected from a certified load interval, but this work does not infer such an interval from traffic. No finite-sample estimation or predictor-robustness statement is implicit in the oracle model. Learning-augmented scheduling results emphasize consistency, robustness, or fairness under explicit prediction errors~\cite{purohit2018,mitzenmacher2022,li2025,cosson2026}; those guarantees cannot be obtained by relabeling a certified interval as a learned estimate. A statistically grounded extension would need a sampling protocol, confidence event, nonstationarity model, and fallback analysis when the interval misses the true load.

Similarly, an unbounded-length code can have finite expected length under a specified prior while still requiring arbitrarily large expanded parameters. Those are distinct resource notions. The speed-augmentation theorem changes actual service capacity and keeps the unit-speed comparator explicit; it is not an alternative interpretation of the unit-speed advice result.

\subsection{Discriminating next steps}\label{sec:next-steps}
The unresolved directions admit concrete tests that distinguish a real extension from a reformulation.  For the equality case, a positive result must exhibit a policy with exact factor $\theta(\rho)$ and prove a tail bound without a strictly stable constant-rate PS comparison.  A natural candidate would let the oldest-job reserve approach $1/\theta$ only when its prefix certificate is endangered, returning more capacity to sharing at other times.  Such a rule is useful only if one can prove both a cumulative prefix-service inequality and a stochastic comparison whose effective spare capacity does not vanish on the trajectories that dominate the tail.  Conversely, an equality lower bound must replace the strict truncation inequality~\eqref{eq:K-choice}; merely taking $K\to\infty$ in the present constants sends the delayed-cohort margin to zero.

For non-Poisson input, three replacements are independently necessary.  A renewal or mixing process must supply a growing-window fluid event after a cycle root; the root and its future input must be coupled in the correct Palm law; and the positive proof needs either an invariant residual description or another conditional response bound for the comparison queue.  Verifying only the first item would leave both the typical-arrival interpretation and the permanent-customer construction unsupported.  A useful first target is a renewal process with a nonlattice interarrival law and finite variance, where regenerative cycles remain available but PASTA does not.

For parallel servers, the first obligation is deterministic rather than asymptotic.  One must identify the correct finite-input maximum-flow optimum when migration, splitting, and preemption are specified.  Total workload divided by total capacity is only a lower bound because jobs cannot necessarily use all servers simultaneously.  An exact small oracle could expose whether an oldest-job guard on one server, a global fractional reserve, or a migratory rule has any bounded ratio before a heavy-tail analysis is attempted.  Only after that benchmark is settled would a multiserver analogue of the delayed-cohort obstruction be meaningful.

\subsection{Evidence boundary}
The exact finite checks exercise the algebraic service accounting, the finite workload oracle, codebook construction, balance recurrences, and selected boundary cases.  They are useful because they expose implementation mistakes such as counting capped work as bounded jobs or using the ordinary geometric count for the permanent-customer system.  They do not prove regular variation, convergence of a stochastic process, or any theorem quantified over all finite inputs.  Those claims rest on the written arguments.

The mathematical conclusion is therefore confined but joint: under one declared observation model, the same policy has both guarantees, and a lower bound applies to every policy in the comparison class. The proof does not classify equality, optimize the leading constant, estimate advice from data, or establish a multiserver analogue. Those limitations mark concrete next problems rather than exceptions hidden inside the theorem statement.
